\section{Machine Learning}

\begin{frame}{O que é Machine Learning?}
  \begin{block}{Uma definição de trabalho}
    Em resumo, um método de Machine Learning pode ser visto como um
    \alert{método numérico na presença de um conjunto de dados} e de um \alert{objetivo}:
    dado um modelo paramétrico $u_\theta$ e uma \emph{função de perda} $\mathcal{L}(\theta)$
    que mede o quão longe estamos do objetivo, ajustamos os parâmetros $\theta$ por
    otimização (tipicamente descida do gradiente) até minimizá-la.
  \end{block}

  \vspace{0.3cm}
  Os três ingredientes que reencontraremos em \emph{todos} os métodos a seguir:
  \begin{enumerate}
    \item um \textbf{modelo} flexível e diferenciável ($u_\theta$);
    \item um \textbf{objetivo} codificado numa perda (dados, física, ou ambos);
    \item um \textbf{otimizador} que segue o gradiente da perda.
  \end{enumerate}
\end{frame}

\begin{frame}{Um breve histórico}
  \begin{columns}[T,totalwidth=\textwidth]
    \begin{column}{0.54\textwidth}
      \begin{itemize}
        \item \textbf{1943} --- neurônio formal de \textcite{mcculloch1943};
        \item \textbf{1958} --- \emph{perceptron} de \textcite{rosenblatt1958}, o primeiro
              modelo treinável;
        \item \textbf{1986} --- \emph{backpropagation} \parencite{rumelhart1986}: torna
              prático o treino de redes profundas;
        \item \textbf{1989} --- teorema da aproximação universal
              \parencite{cybenko1989,hornik1989};
        \item \textbf{2015+} --- \emph{deep learning} \parencite{lecun2015deep} e a
              revolução de IA em curso.
      \end{itemize}
    \end{column}
    \begin{column}{0.46\textwidth}
      \begin{information}[Por que agora?]
        Vivemos uma \alert{expansão sem precedentes} da área: \emph{hardware} (GPUs),
        \emph{frameworks} de diferenciação automática (PyTorch) e uma vasta literatura
        tornaram esses métodos acessíveis.
      \end{information}
      \vspace{0.2cm}
      \begin{remark}
        Um dos objetivos deste trabalho é justamente \textbf{tirar proveito} desse
        ferramental moderno, trazendo-o para a solução numérica de EDPs.
      \end{remark}
    \end{column}
  \end{columns}
\end{frame}

\begin{frame}{O cenário do ML para EDPs}
  Para resolver EDPs, três famílias guiam este trabalho:
  \begin{columns}[T,totalwidth=\textwidth]
    \begin{column}{0.33\textwidth}
      \begin{block}{PINN}
        \footnotesize Rede neural com a \textbf{física na perda}
        \parencite{raissi2019}. Resolve \emph{uma} instância da EDP.
      \end{block}
    \end{column}
    \begin{column}{0.33\textwidth}
      \begin{block}{FNO}
        \footnotesize Aprende o \textbf{operador solução} entre espaços de funções, de
        forma puramente \emph{data-driven} \parencite{li2020fourier}.
      \end{block}
    \end{column}
    \begin{column}{0.33\textwidth}
      \begin{block}{PINO}
        \footnotesize Operador neural com \textbf{física na perda}
        \parencite{li2023pino}: o híbrido dos dois.
      \end{block}
    \end{column}
  \end{columns}

  \vspace{0.4cm}
  \begin{information}[Lacuna que motiva o trabalho]
    Métodos de Machine Learning já foram aplicados a diversas equações de águas rasas e ondas
    dispersivas, mas, até onde sabemos, \alert{nenhum estudo avaliou sistematicamente FNO,
    PINO e PINN no sistema de Boussinesq paramétrico 1D}, com não linearidade e dispersão
    variando simultaneamente.
  \end{information}
\end{frame}
